% \thispagestyle{empty}
% \begin{figure}[H]
%     \begin{center}
%         \includegraphics[width=1.05\linewidth]{Thuc_Hanh/Danhgia/DanhGia-part-1.pdf}
%      \end{center}
%  \end{figure}
%  \newpage
% \thispagestyle{empty}
% \begin{figure}[H]
%     \begin{center}
%         \includegraphics[width=1.05\linewidth]{Thuc_Hanh/Danhgia/DanhGia-part-3.pdf}
%      \end{center}
%  \end{figure}
% \newpage
% \thispagestyle{empty}
% \begin{figure}[H]
%     \begin{center}
%         \includegraphics[width=1.05\linewidth]{Thuc_Hanh/Danhgia/DanhGia-part-5.pdf}
%      \end{center}
%  \end{figure}
% \cleardoublepage

% \includepdf[pages=-]{Thuc_Hanh/Danhgia/DanhGia.pdf}

% ===================================================
%  MACRO: in một phiếu đánh giá CBHD (dùng lại 2 lần)
% ===================================================
\newcommand{\phieuCBHD}{%
  \begingroup
  \newgeometry{top=1.4cm,bottom=1.4cm,left=1.8cm,right=1.5cm}
  \thispagestyle{empty}
  \setlength{\parindent}{0pt}
  \setlength{\parskip}{2pt}
  \renewcommand{\baselinestretch}{1.1}\selectfont

  % ---------- Tiêu đề ----------
  \begin{center}
    {\fontsize{12pt}{14pt}\selectfont ĐẠI HỌC BÁCH KHOA HÀ NỘI}\\[1pt]
    {\fontsize{12pt}{14pt}\selectfont\bfseries\underline{TRƯỜNG ĐIỆN -- ĐIỆN TỬ}}\\[7pt]
    {\fontsize{13pt}{15pt}\selectfont\bfseries ĐÁNH GIÁ ĐỒ ÁN TỐT NGHIỆP}\\[1pt]
    {\fontsize{12pt}{14pt}\selectfont\bfseries (DÀNH CHO CÁN BỘ HƯỚNG DẪN)}
  \end{center}

% ---------- Thông tin chung ----------
\vspace{-10pt}
\begin{table}[H]
\centering
\renewcommand{\arraystretch}{1.14}
\begin{tabular}{p{3.4cm} p{13.5cm}}
Tên đề tài: & {\bfseries  Nghiên cứu thiết kế và chế tạo hệ thống IoT giám sát vi khí hậu chuồng nuôi gia cầm ứng dụng LoRa và Điện toán biên} \\
Họ tên SV: & {\bfseries Nguyễn Tiến Đạt} \hfill MSSV: {\bfseries 20223689} \\
Cán bộ hướng dẫn: & {\bfseries PGS.TS. Trần Quang Vinh} \\
\end{tabular}
\end{table}
\vspace{-6pt}
%----------------------------------------

  \footnotesize
  \renewcommand{\arraystretch}{1.18}
  \setlength{\tabcolsep}{4pt}

  \noindent
  \begin{tabular}{|>{\centering\arraybackslash\bfseries}m{0.62cm}%
                  |>{\centering\arraybackslash}m{2.05cm}%
                  |m{12.55cm}%
                  |>{\centering\arraybackslash}m{1.3cm}|}
  \hline
  \textbf{STT} & \textbf{Tiêu chí (Điểm tối đa)} & \centering\arraybackslash\textbf{Hướng dẫn đánh giá tiêu chí} & \textbf{Điểm tiêu chí} \\
  \hline

  1 &
  \textbf{Thái độ làm việc} \newline\textbf{(2,5 điểm)} &
  Nghiêm túc, tích cực và chủ động trong quá trình làm ĐATN.

  Hoàn thành đầy đủ và đúng tiến độ các nội dung được GVHD giao.
  & 2,5\\
  \hline

  2 &
  \textbf{Kỹ năng viết quyển ĐATN} \newline\textbf{(2 điểm)} &
  Trình bày đúng mẫu quy định, bố cục các chương logic và hợp lý: Bảng biểu, hình
  ảnh rõ ràng, có tiêu đề, được đánh số thứ tự và được giải thích hay để cập đến
  trong đồ án, có căn lề, dấu cách sau dấu chấm, dấu phẩy, có mở đầu chương và
  kết luận chương, có liệt kê tài liệu tham khảo và có trích dẫn, v.v.

  Kỹ năng diễn đạt, phân tích, giải thích, lập luận: Cấu trúc câu rõ ràng, văn
  phong khoa học, lập luận logic và có cơ sở, thuật ngữ chuyên ngành phù hợp, v.v.
  & 2\\
  \hline

  3 &
  \textbf{Nội dung và kết quả đạt được} \newline\textbf{(5 điểm)} &
  Nêu rõ tính cấp thiết, ý nghĩa khoa học và thực tiễn của đề tài, các vấn đề và
  các giả thuyết, phạm vi ứng dụng của đề tài. Thực hiện đầy đủ quy trình nghiên
  cứu: Đặt vấn đề, mục tiêu đề ra, phương pháp nghiên cứu/ giải quyết vấn đề, kết
  quả đạt được, đánh giá và kết luận.

  Nội dung và kết quả được trình bày một cách logic và hợp lý, được phân tích và
  đánh giá thỏa đáng. Biện luận phân tích kết quả mô phỏng/ phần mềm/ thực
  nghiệm, so sánh kết quả đạt được với kết quả trước đó có liên quan.

  Chỉ rõ phù hợp giữa kết quả đạt được và mục tiêu ban đầu đề ra đồng thời cung
  cấp lập luận để đề xuất hướng giải quyết có thể thực hiện trong tương lai. Hàm
  lượng khoa học/ độ phức tạp cao, có tính mới/tính sáng tạo trong nội dung và
  kết quả của đồ án.
  & 5\\
  \hline

  4 &
  \textbf{Điểm thành tích} \newline\textbf{(1 điểm)} &
  Có bài báo KH được đăng hoặc chấp nhận đăng/ đạt giải SV NCKH giải 3 cấp
  Trường trở lên/ Các giải thưởng khoa học trong nước, quốc tế từ giải 3 trở lên/
  Có đăng ký bằng phát minh sáng chế. \textbf{(1 điểm)}

  Được báo cáo tại hội đồng cấp Trường trong hội nghị SV NCKH nhưng không đạt
  giải từ giải 3 trở lên/ Đạt giải khuyến khích trong cuộc thi khoa học trong nước,
  quốc tế/ Kết quả đồ án là sản phẩm ứng dụng có tính hoàn thiện cao, yêu cầu
  khối lượng thực hiện lớn. \textbf{(0,5 điểm)}
  & 0\\
  \hline

  \multicolumn{3}{|r|}{\textbf{Điểm tổng các tiêu chí:}} & 9,5\\
  \hline
  \multicolumn{3}{|r|}{\textbf{Điểm hướng dẫn:}} & 9,5\\
  \hline
  \end{tabular}

  % ---------- Chữ ký ----------
  %\vspace{6pt}
  \begin{flushright}
    \fontsize{12pt}{20pt}\selectfont\textit{Cán bộ hướng dẫn}\\[22pt]
    \fontsize{12pt}{24pt}\selectfont\textbf{PGS.TS. Trần Quang Vinh}
  \end{flushright}

  \noindent{\fontsize{9.5pt}{11pt}\selectfont\textit{%
    Điểm từng tiêu chí cho lẻ đến 0,5.
    Nếu Điểm tổng các tiêu chí $>$ 10 thì Điểm hướng dẫn làm tròn thành 10.}}

  \restoregeometry
  \renewcommand{\baselinestretch}{1.2}\selectfont
  \setlength{\parskip}{6pt}
  \setlength{\parindent}{1cm}
  \endgroup
}% end \phieuCBHD

% ===================================================
%  MACRO: phiếu đánh giá CBPB (dùng lại 2 lần)
% ===================================================
\newcommand{\phieuCBPB}{%
  \begingroup
  \newgeometry{top=1.4cm,bottom=1.4cm,left=1.8cm,right=1.5cm}
  \thispagestyle{empty}
  \setlength{\parindent}{0pt}
  \setlength{\parskip}{2pt}
  \renewcommand{\baselinestretch}{1.1}\selectfont

  % ---------- Tiêu đề ----------
  \begin{center}
    {\fontsize{12pt}{14pt}\selectfont ĐẠI HỌC BÁCH KHOA HÀ NỘI}\\[1pt]
    {\fontsize{12pt}{14pt}\selectfont\bfseries\underline{TRƯỜNG ĐIỆN -- ĐIỆN TỬ}}\\[7pt]
    {\fontsize{13pt}{15pt}\selectfont\bfseries ĐÁNH GIÁ ĐỒ ÁN TỐT NGHIỆP}\\[1pt]
    {\fontsize{12pt}{14pt}\selectfont\bfseries (DÀNH CHO CÁN BỘ PHẢN BIỆN)}
  \end{center}

 % ---------- Thông tin chung ----------
\vspace{-10pt}
\begin{table}[H]
\centering
\renewcommand{\arraystretch}{1.14}
\begin{tabular}{p{3.4cm} p{13.5cm}}
Tên đề tài: & {\bfseries  Nghiên cứu thiết kế và chế tạo hệ thống IoT giám sát vi khí hậu chuồng nuôi gia cầm ứng dụng LoRa và Điện toán biên} \\
Họ tên SV: & {\bfseries Nguyễn Tiến Đạt} \hfill MSSV: {\bfseries 20223689} \\
Cán bộ phản biện: & ................................................................................................................ \\
\end{tabular}
\end{table}
\vspace{-6pt}
%----------------------------------------

  % ---------- Bảng đánh giá ----------
  \footnotesize
  \renewcommand{\arraystretch}{1.18}
  \setlength{\tabcolsep}{4pt}

  \noindent
  \begin{tabular}{|>{\centering\arraybackslash\bfseries}m{0.62cm}%
                  |>{\centering\arraybackslash}m{2.05cm}%
                  |m{12.55cm}%
                  |>{\centering\arraybackslash}m{1.3cm}|}
  \hline
  \textbf{STT} & \textbf{Tiêu chí (Điểm tối đa)} & \centering\arraybackslash\textbf{Hướng dẫn đánh giá tiêu chí} & \textbf{Điểm tiêu chí} \\
  \hline

  1 &
  \textbf{Trình bày quyển ĐATN} \newline\textbf{(4 điểm)} &
  Trình bày đúng mẫu quy định, bố cục các chương logic và hợp lý: Bảng biểu, hình
  ảnh rõ ràng, có tiêu đề, được đánh số thứ tự và được giải thích hay để cập đến
  trong đồ án, có căn lề, dấu cách sau dấu chấm, dấu phẩy, có mở đầu chương và
  kết luận chương, có liệt kê tài liệu tham khảo và có trích dẫn, v.v.

  Kỹ năng diễn đạt, phân tích, giải thích, lập luận: Cấu trúc câu rõ ràng, văn
  phong khoa học, lập luận logic và có cơ sở, thuật ngữ chuyên ngành phù hợp, v.v.
  & \\
  \hline

  2 &
  \textbf{Nội dung và kết quả đạt được} \newline\textbf{(5,5 điểm)} &
  Nêu rõ tính cấp thiết, ý nghĩa khoa học và thực tiễn của đề tài, các vấn đề và
  các giả thuyết, phạm vi ứng dụng của đề tài. Thực hiện đầy đủ quy trình nghiên
  cứu: Đặt vấn đề, mục tiêu đề ra, phương pháp nghiên cứu/ giải quyết vấn đề, kết
  quả đạt được, đánh giá và kết luận.

  Nội dung và kết quả được trình bày một cách logic và hợp lý, được phân tích và
  đánh giá thỏa đáng. Biện luận phân tích kết quả mô phỏng/ phần mềm/ thực
  nghiệm, so sánh kết quả đạt được với kết quả trước đó có liên quan.

  Chỉ rõ phù hợp giữa kết quả đạt được và mục tiêu ban đầu đề ra đồng thời cung
  cấp lập luận để đề xuất hướng giải quyết có thể thực hiện trong tương lai. Hàm
  lượng khoa học/ độ phức tạp cao, có tính mới/tính sáng tạo trong nội dung và
  kết quả của đồ án.
  & \\
  \hline

  3 &
  \textbf{Điểm thành tích} \newline\textbf{(1 điểm)} &
  Có bài báo KH được đăng hoặc chấp nhận đăng/ đạt giải SV NCKH giải 3 cấp
  Trường trở lên/ Các giải thưởng KH quốc tế/ trong nước từ giải 3 trở lên/
  Có đăng ký bằng phát minh sáng chế. \textbf{(1 điểm)}

  Được báo cáo tại hội đồng cấp Trường trong hội nghị SV NCKH nhưng không đạt
  giải từ giải 3 trở lên/ Đạt giải khuyến khích trong các kỳ thi quốc gia và
  quốc tế khác về chuyên ngành. \textbf{(0,5 điểm)}
  & \\
  \hline

  \multicolumn{3}{|r|}{\textbf{Điểm tổng các tiêu chí:}} & \\
  \hline
  \multicolumn{3}{|r|}{\textbf{Điểm phản biện:}} & \\
  \hline
  \end{tabular}

  % ---------- Chữ ký ----------
  \vspace{8pt}
  \begin{flushright}
    \fontsize{12pt}{14pt}\selectfont\textit{Cán bộ phản biện}\\[2pt]
    \fontsize{11pt}{13pt}\selectfont\textit{(Ký và ghi rõ họ tên)}\\[24pt]
    \phantom{PGS.TS. Trần Quang Vinh}
  \end{flushright}

  \vspace{4pt}
  \noindent{\fontsize{9.5pt}{11pt}\selectfont\textit{%
    Điểm từng tiêu chí cho lẻ đến 0,5.
    Nếu Điểm tổng các tiêu chí $>$ 10 thì Điểm hướng dẫn làm tròn thành 10.}}

  \restoregeometry
  \renewcommand{\baselinestretch}{1.2}\selectfont
  \setlength{\parskip}{6pt}
  \setlength{\parindent}{1cm}
  \endgroup
}% end \phieuCBPB

% ===================================================
%  In hai bản phiếu CBHD + hai bản phiếu CBPB
% ===================================================
\phieuCBHD
\cleardoublepage

\phieuCBPB
\cleardoublepage